\begin{abstract}
Intracranial hemorrhage (ICH) is a life-threatening neurological emergency
in which blood accumulates within the cranial vault following traumatic
injury or spontaneous vascular rupture.
Rapid and accurate identification of the hemorrhage type on non-contrast
head computed tomography (CT) is critical for treatment triage, yet
radiologist availability and inter-reader variability remain bottlenecks
in emergency settings.
This paper presents a two-stage deep learning framework for the automatic
detection and subtype classification of ICH across six categories:
no hemorrhage, epidural, intraparenchymal, intraventricular, subarachnoid,
and subdural.
In Stage~1, an ensemble of three convolutional neural networks---DenseNet-121,
DenseNet-169, and SE-ResNeXt-101---processes adjacent CT slices encoded as
three-channel images after Hounsfield Unit windowing.
The ensemble logits are averaged to produce robust per-slice predictions.
In Stage~2, a bidirectional gated recurrent unit network with a
parallel one-dimensional convolutional branch refines the per-slice
predictions using inter-slice context within each CT study, additionally
incorporating slice thickness as a metadata feature derived from the DICOM
header.
The complete system is trained and evaluated on the public CQ500 dataset
comprising 491 CT studies with consensus radiologist labels.
Experimental results demonstrate competitive performance across all six
hemorrhage subtypes as measured by area under the receiver operating
characteristic curve, sensitivity, and specificity with 95\% confidence
intervals estimated via bootstrap resampling.
\end{abstract}

\begin{IEEEkeywords}
Intracranial hemorrhage, CT imaging, deep learning, DenseNet, SE-ResNeXt,
bidirectional GRU, multi-label classification, RSNA 2019, CQ500
\end{IEEEkeywords}
